\subsection{Berechnung der Ausgangsspannung einer zweistufigen OPV-Schaltung}
% Alt: Aufgabe 6
\Aufgabe{
    \label{Aufg6}
    Berechnen Sie $U_\mathrm{A}$ mit den folgenden Werten: \newline
    $R_1 = \mathrm{10\,k\Omega}$, $R_2 = \mathrm{20\,k\Omega}$, $R_3 = \mathrm{30\,k\Omega}$, $R_4 = \mathrm{10\,k\Omega}$, $R_5 = \mathrm{20\,k\Omega}$

    \begin{figure}[H]
        \centering
        \resizebox{0.95\textwidth}{!}{\input{Tikz/src/FigOPVzweistufig.tex}\alttext{Das Bild zeigt einen elektronischen Schaltkreis mit sechs Widerständen und zwei Operationsverstärkern. Links sind drei Spannungsquellen mit den Bezeichnungen \(U_{\mathrm{E1}}\), \(U_{\mathrm{E2}}\) und \(U_{\mathrm{E3}}\), die von oben nach unten angeordnet sind. Die Widerstände sind beschriftet als \(R_1\), \(R_2\), \(R_3\), \(R_4\), \(R_5\) und \(R_6\).
                \(R_1\) ist oben links parallel zu \(R_2\) darunter geschaltet. Darunter befindet sich \(R_3\), der mit einem Knotenpunkt verbunden ist. Von diesem Knoten führt eine Leitung zu dem ersten Operationsverstärker N1. Zwischen dem Eingang des Operationsverstärkers N1 und seinem Ausgang liegt der Widerstand \(R_4\) in einer Rückkopplungsschleife. Der Ausgang von N1 ist über den Widerstand \(R_5\) mit dem Eingang des zweiten Operationsverstärkers N2 verbunden.
                Zwischen dem Eingang und Ausgang von N2 befindet sich der Widerstand \(R_6\) in einer Rückkopplungsschleife. Der Ausgang des zweiten Operationsverstärkers ist als Spannung \(U_\mathrm{A}\) bezeichnet. Alle Verbindungen sind durch Leitungen dargestellt, die die Bauteile miteinander verbinden und auf Masse bezogen sind.}}
        \label{fig:FigOPVzweistufig}
    \end{figure}
}

\Loesung{
    Formel:
    \begin{align*}
        U_\mathrm{N1} = - \left( \frac{R_\mathrm{N}}{R_\mathrm{1}} U_\mathrm{E1} + \frac{R_\mathrm{N}}{R_2} U_\mathrm{E2} + \frac{R_\mathrm{N}}{R_\mathrm{3}} U_\mathrm{E3} \right)
    \end{align*}
    \begin{align*}
        U_\mathrm{A} = -\frac{R_\mathrm{6}}{R_\mathrm{5}} U_\mathrm{N1}
    \end{align*}

    Berechnung:

    Berechnung von $U_\mathrm{N1}$
    \begin{align*}
        U_\mathrm{N1} & = - \left( \frac{10\,\mathrm{k}\Omega}{10\,\mathrm{k}\Omega} U_\mathrm{E1} + \frac{10\,\mathrm{k}\Omega}{20\,\mathrm{k}\Omega} U_\mathrm{E2} + \frac{10\,\mathrm{k}\Omega}{30\,\mathrm{k}\Omega} U_\mathrm{E3} \right) \\
                      & = - \left( U_\mathrm{E1} + 0,5 U_\mathrm{E2} + \frac{1}{3} U_\mathrm{E3} \right)
    \end{align*}

    Berechnung von $U_\mathrm{A}$
    \begin{align*}
        U_\mathrm{A} & = -2 \cdot U_\mathrm{N1}                                                                           \\
                     & = -2 \left( - \left( U_\mathrm{E1} + 0,5 U_\mathrm{E2} + \frac{1}{3} U_\mathrm{E3} \right) \right) \\
                     & = 2 \left( U_\mathrm{E1} + 0,5 U_\mathrm{E2} + \frac{1}{3} U_\mathrm{E3} \right)
    \end{align*}

    Die Ausgangsspannung $U_\mathrm{A}$ beträgt:
    \begin{align*}
        U_\mathrm{A} & = 2 U_\mathrm{E1} + U_\mathrm{E2} + \frac{2}{3} U_\mathrm{E3}
    \end{align*}

    Siehe: Abschnitt \ref{fig: Verschaltung Verstaerker} und Tabelle \ref{tab:Grundschaltungen}
}